\begin{center}
\LARGE
Techniques for the symbolic analysis of distortion and intermodulation

\end{center}

\begin{center}
\Large
Georges Gielen,  Piet Wambacq, Willy Sansen
\end{center}

\noindent
Date:  July 19th (Friday) \\
Time: 09:40-10:00\\

\begin{center}
\large
Abstract
\end{center}

The booming market of wireless communications demands for integrated
radio-frequency (RF) circuits. The design of these RF circuits is
dominated by second-order effects such as noise figure, harmonic 
distortion and intermodulation products. Symbolic simulation tools
allow designers to analyze these effects symbolically or mixed
symbolic-numerically, which offers them a better understanding in the
mechanisms of these effects and therefore also provides them indications
on how to improve the quality of their designs.

In this presentation, techniques will be presented for the symbolic
analysis of harmonic distortion and intermodulation in weakly nonlinear
analog integrated circuits. Examples will be presented for practical
circuits such as opamps, mixers, etc.



